\documentclass{article}
\usepackage{lmodern}
\usepackage{amsmath}
\usepackage{listings}
\usepackage{fullpage}
\usepackage{parskip}
\usepackage{xcolor}
\lstset{
	basicstyle=\ttfamily,
	columns=fullflexible,
	frame=single,
	breaklines=true,
	postbreak=\mbox{\textcolor{red}{$\hookrightarrow$}\space},
}
\begin{document}
\title{Experiment `p\_6\_6\_alternating\_sum\_seq' Results}
\author{\today}
\date{}
\maketitle
\textbf{Experiment outcome: }SUCCESS
\\\textbf{Bad responses: }0
\\\textbf{Responses containing}~\texttt{assume}~\textbf{: }0
\\\textbf{Responses containing}~\texttt{expect}~\textbf{: }0
\\\textbf{Resolution attempts: }1
\\\textbf{Hard fails (resolution): }0
\\\textbf{Soft fails (resolution): }0
\\\textbf{Verification attempts: }1
\section*{Problem Specification}
\textbf{Problem name: }p\_6\_6\_alternating\_sum\_seq
\\\textbf{Natural language statement: }Compute the alternating sum of all elements in an array. For example, if your method reads the input [1, 4, 9, 16, 9, 7, 4, 9, 11] then it computes 1 – 4 + 9 – 16 + 9 – 7 + 4 – 9 + 11 = –2.
\\\textbf{Method signature: }p\_6\_6\_alternating\_sum\_seq(inputs: seq<int>) returns (alt\_sum: int)
\subsection*{Ensures}
\begin{itemize}
\item \texttt{alt\_sum == sum\_array(inputs, 1)}
\end{itemize}
\subsection*{Functional Code Given}
\begin{lstlisting}
function sum_array(inputs: seq<int>, i: int): int
{
  if |inputs| == 0 then 0
  else inputs[0] * i + sum_array(inputs[1..], i * -1)
}
\end{lstlisting}
\clearpage
\section*{GenAI interactions}
Below you will find all interactions between the `user' (program) and the `assistant' (OpenAI).
\subsection*{Program $\rightarrow$ GenAI}
\begin{lstlisting}
You are given the following task to perform in Dafny:

Compute the alternating sum of all elements in an array. For example, if your method reads the input [1, 4, 9, 16, 9, 7, 4, 9, 11] then it computes 1 – 4 + 9 – 16 + 9 – 7 + 4 – 9 + 11 = –2.

The signature should be:

p_6_6_alternating_sum_seq(inputs: seq<int>) returns (alt_sum: int)

The method should respect the following contract:

ensures alt_sum == sum_array(inputs, 1)

The contract uses the following dafny code:

function sum_array(inputs: seq<int>, i: int): int
{
  if |inputs| == 0 then 0
  else inputs[0] * i + sum_array(inputs[1..], i * -1)
}

These function(s) must not be used in your implementation of the method.

Produce and show only the Dafny body of this method, including the curly braces that surround it. Do not show the signature nor contract. You must not use 'assume' nor 'expect' anywhere in your code.
\end{lstlisting}
\subsection*{GenAI $\rightarrow$ Program}
\textbf{System fingerprint: }null
\\\textbf{ID: }chatcmpl-Dn1ZODDV1pdjTQpYK6xpRFWggu34s
\\\textbf{Model: }gpt-5.5-2026-04-23
\\\textbf{Created at: }1780576478
\\\textbf{Finish reason: }stop
\\\textbf{Usage: }550 tokens in, and 6386 tokens out
\begin{lstlisting}
{
  if |inputs| == 0 {
    alt_sum := 0;
  } else if |inputs| == 1 {
    assert inputs[1..] == [];
    alt_sum := inputs[0];
  } else {
    assert 2 <= |inputs|;
    assert 0 < |inputs[1..]|;
    assert (inputs[1..])[0] == inputs[1];
    assert (inputs[1..])[1..] == inputs[2..];
    var rest := p_6_6_alternating_sum_seq(inputs[2..]);
    alt_sum := inputs[0] - inputs[1] + rest;
  }
}
\end{lstlisting}
\section*{Final Program}
\begin{lstlisting}
function sum_array(inputs: seq<int>, i: int): int
{
  if |inputs| == 0 then 0
  else inputs[0] * i + sum_array(inputs[1..], i * -1)
}

method p_6_6_alternating_sum_seq(inputs: seq<int>) returns (alt_sum: int)
	ensures alt_sum == sum_array(inputs, 1)
{
  if |inputs| == 0 {
    alt_sum := 0;
  } else if |inputs| == 1 {
    assert inputs[1..] == [];
    alt_sum := inputs[0];
  } else {
    assert 2 <= |inputs|;
    assert 0 < |inputs[1..]|;
    assert (inputs[1..])[0] == inputs[1];
    assert (inputs[1..])[1..] == inputs[2..];
    var rest := p_6_6_alternating_sum_seq(inputs[2..]);
    alt_sum := inputs[0] - inputs[1] + rest;
  }
}
\end{lstlisting}
\section*{Total Token Usage}
\textbf{Input tokens: }550
\\\textbf{Output tokens: }6386
\\\textbf{Reasoning tokens: }6145
\\\textbf{Sum of `total tokens': }6936
\section*{Experiment Timings}
\textbf{Overall Experiment} started at 1780576477582, ended at 1780576651446, lasting 173864ms (173.86 seconds)
\\\textbf{Iteration \#1} started at 1780576477641, ended at 1780576651446, lasting 173805ms (173.81 seconds)
\end{document}
